\par\Needspace{6\baselineskip}\medskip
\SourceLeaf{149}\pdfbookmark[2]{107.}{article-107}
\begin{center}107.\end{center}
Bezeichnen wir mit $\chi$ einen beliebigen Bogen zwischen $-180^0$ und $+180^0$ und mit $s$ den Sinus des Bogens $\tfrac12\chi$, so gilt bekanntlich
\[
\tfrac12\chi=s+\tfrac13\cdot\tfrac12s^3+\tfrac15\cdot\frac{1\mathpunct{.}3}{2\mathpunct{.}4}s^5+\tfrac17\cdot\frac{1\mathpunct{.}3\mathpunct{.}5}{2\mathpunct{.}4\mathpunct{.}6}s^7+\text{usw.}
\]
\noindent Ferner gilt
\[
\tfrac12\sin\chi=s\BareRoot{(1-ss)}=s-\tfrac12s^3-\frac{1\mathpunct{.}1}{2\mathpunct{.}4}s^5-\frac{1\mathpunct{.}1\mathpunct{.}3}{2\mathpunct{.}4\mathpunct{.}6}s^7-\text{usw.}
\]
\noindent und daher
\[
\begin{aligned}
\chi-\sin\chi=4\Bigl(&\tfrac13s^3+\tfrac15\cdot\tfrac12s^5\\[.5em]
&+\tfrac17\cdot\frac{1\mathpunct{.}3}{2\mathpunct{.}4}s^7+\tfrac19\cdot\frac{1\mathpunct{.}3\mathpunct{.}5}{2\mathpunct{.}4\mathpunct{.}6}s^9+\text{usw.}\Bigr).
\end{aligned}
\]
\noindent Wir setzen in dieser Reihe für $s$ nacheinander $\tfrac12\BareRoot{\dfrac{r+r'-\rho}{a}}$ und $\tfrac12\BareRoot{\dfrac{r+r'+\rho}{a}}$ ein und multiplizieren die Ergebnisse mit $a^{\frac32}$. Daraus ergeben sich der Reihe nach die Reihen
\[
\begin{aligned}
&\tfrac16(r+r'-\rho)^{\frac32}+\tfrac1{80}\cdot\frac1a(r+r'-\rho)^{\frac52}\\[.5em]
&\qquad+\tfrac3{1792}\cdot\frac1{aa}(r+r'-\rho)^{\frac72}+\tfrac5{18432}\cdot\frac1{a^3}(r+r'-\rho)^{\frac92}+\text{usw}
\end{aligned}
\]
\[
\begin{aligned}
&\tfrac16(r+r'+\rho)^{\frac32}+\tfrac1{80}\cdot\frac1a(r+r'+\rho)^{\frac52}\\[.5em]
&\qquad+\tfrac3{1792}\cdot\frac1{aa}(r+r'+\rho)^{\frac72}+\tfrac5{18432}\cdot\frac1{a^3}(r+r'+\rho)^{\frac92}+\text{usw.,}
\end{aligned}
\]
\noindent deren Summen wir mit $T$ und $U$ bezeichnen. Da $2\sin\tfrac12\delta=\pm\BareRoot{\dfrac{r+r'-\rho}{a}}$ ist, wobei das obere oder untere Vorzeichen gilt, je nachdem, ob $2f$ unter oder über $180^0$ liegt, ergibt sich unmittelbar $a^{\frac32}(\delta-\sin\delta)=\pm T$ mit derselben Wahl des Vorzeichens. Nimmt man für $\epsilon$ den kleineren Wert unter $180^0$, so gilt entsprechend $a^{\frac32}(\epsilon-\sin\epsilon)=U$. Nimmt man den anderen Wert, der den ersten zu $360^0$ ergänzt, so ergibt sich offenbar $a^{\frac32}(\epsilon-\sin\epsilon)=a^{\frac32}\mathpunct{.}360^0-U$. Daraus ergeben sich zwei Werte für die Zeit $t$
\[
\frac{U\mp T}{k}\qquad\text{und}\qquad\frac{a^{\frac32}\mathpunct{.}360^0}{k}-\frac{U\pm T}{k}\mathpunct{\cdot}
\]
